\documentclass[12pt,a4paper]{article}
\usepackage[utf8]{inputenc}
\usepackage[T1]{fontenc}
\usepackage{thaisup}
\usepackage{graphicx}
\usepackage{tikz}
\usepackage{amsmath,amssymb}
\usepackage{geometry}
\geometry{margin=2.5cm}
\usepackage{enumitem}
\usepackage{fancyhdr}
\usepackage{booktabs}
\usepackage{hyperref}

% Header/Footer
\pagestyle{fancy}
\fancyhf{}
\rhead{Part 1}
\lhead{ครั้งที่ 1}
\rfoot{หน้า \thepage}

% Question formatting
\setlist[enumerate]{fullwidth, itemindent=2em, label=\arabic*.}
\renewenvironment{enumerate}{\begin{trivlist}\item\薄}{\end{trivlist}}
\let\薄\item

% Line separator
\newcommand{\HRule}{\\ \hline \\}

% Answer box
\\newcounter{savedenum}
\usepackage{etoolbox}
\AtBeginEnvironment{enumerate}{\\stepcounter{savedenum}}

\begin{document}

\begin{document}
\title{\textbf{Part 1: ครั้งที่ 1}}
\subtitle{บทนำ: หน่วย มาตรฐาน และเลขนัยสำคัญ}
\date{\today}
\maketitle
\section*{หัวข้อที่เรียน}
\begin{itemize}
\item หน่วย
\item มาตรฐาน
\item เลขนัยสำคัญ
\item ความไม่แน่นอนในการวัด
\end{itemize}
\newpage
\section*{แบบฝึกหัดในคาบ (25 ข้อ)}
\begin{enumerate}
\item[1] ข้อใดต่อไปนี้ไม่ใช่หน่วยฐานในระบบ SI?
\begin{enumerate}
\item แอมแปร์ (A)
\item เคลวิน (K)
\item จูล (J)
\item โมล (mol)
\end{enumerate}
\textbf{คำตอบ: } ค
\textit{เฉลย: } หน่วยฐาน SI มี 7 หน่วย ได้แก่ เมตร (m), กิโลกรัม (kg), วินาที (s), แอมแปร์ (A), เคลวิน (K), แคนเดลา (cd), โมล (mol) ส่วนจูล (J) เป็นหน่วยอนุพันธ์
\HRule
\item[2] ค่า $5 \times 10^{-3}$ เมตร มีค่าเท่ากับข้อใด?
\begin{enumerate}
\item 5 มิลลิเมตร
\item 5 ไมโครเมตร
\item 5 เซนติเมตร
\item 5 นาโนเมตร
\end{enumerate}
\textbf{คำตอบ: } ก
\textit{เฉลย: } $5 \times 10^{-3}$ m = $5 \times 10^{-3}$ m = 5 mm (มิลลิ) เพราะ $10^{-3}$ = มิลลิ (m)
\HRule
\item[3] ปริมาณ $3.5$ กิโลกรัม มีค่าเท่ากับกี่กรัม?
\begin{enumerate}
\item 3.5 g
\item 35 g
\item 350 g
\item 3500 g
\end{enumerate}
\textbf{คำตอบ: } ง
\textit{เฉลย: } 1 kg = 1,000 g ดังนั้น $3.5$ kg = $3.5 \times 1000$ g = 3,500 g
\HRule
\item[4] จำนวน $0.00250$ มีเลขนัยสำคัญกี่ตัว?
\begin{enumerate}
\item 2 ตัว
\item 3 ตัว
\item 4 ตัว
\item 5 ตัว
\end{enumerate}
\textbf{คำตอบ: } ข
\textit{เฉลย: } เลขนัยสำคัญคือตัวเลขทั้งหมดที่ไม่ใช่ศูนย์นำหน้า ดังนั้น $0.00250$ มีเลขนัยสำคัญ 3 ตัว คือ 2, 5, 0 (ศูนย์ท้ายมีนัยสำคัญเพราะบอกความละเอียด)
\HRule
\item[5] ข้อใดมีเลขนัยสำคัญ 4 ตัว?
\begin{enumerate}
\item 0.000560
\item 5600
\item 5.600 \times 10^{-3}
\item 560.
\end{enumerate}
\textbf{คำตอบ: } ค
\textit{เฉลย: } ตัวเลือก ก มี 3 ตัว (5, 6, 0), ข มี 2 ตัว (5, 6 ถ้าไม่มีจุดทศนิยม), ง มี 3 ตัว (5, 6, 0) แต่ ค $5.600 \times 10^{-3}$ มี 4 ตัว (5, 6, 0, 0 ศูนย์ท้ายในทศนิยมมีนัยสำคัญ)
\HRule
\item[6] จงหาผลลัพธ์ของ $12.3 + 0.456$
(ให้ตอบโดยใช้เลขนัยสำคัญที่ถูกต้อง)
\textbf{คำตอบ: } 12.76
\textit{วิธีทำ: } การบวกลบใช้การปัดเศษตามตำแหน่งทศนิยมน้อยที่สุด
$12.3$ มีทศนิยม 1 ตำแหน่ง, $0.456$ มี 3 ตำแหน่ง
ดังนั้นตอบที่ 2 ตำแหน่งทศนิยม: $12.3 + 0.456 = 12.756 \approx 12.76$
\HRule
\item[7] จงหาผลลัพธ์ของ $2.5 \times 3.42$
(ให้ตอบโดยใช้เลขนัยสำคัญที่ถูกต้อง)
\textbf{คำตอบ: } 8.6
\textit{วิธีทำ: } การคูณหารใช้เลขนัยสำคัญน้อยที่สุด
$2.5$ มี 2 ตัว, $3.42$ มี 3 ตัว
ดังนั้นตอบที่ 2 ตัว: $2.5 \times 3.42 = 8.55 \approx 8.6$
\HRule
\item[8] รถยนต์วิ่งด้วยความเร็ว $72$ กิโลเมตร/ชั่วโมง มีความเร็วเท่ากับกี่เมตร/วินาที?
\begin{enumerate}
\item 10 m/s
\item 20 m/s
\item 25 m/s
\item 36 m/s
\end{enumerate}
\textbf{คำตอบ: } ข
\textit{เฉลย: } วิธีทำ: $72$ km/h = $72 \times \frac{1000}{3600}$ m/s = $72 \times \frac{5}{18}$ m/s = $20$ m/s
\HRule
\item[9] น้ำ $1$ ลูกบาศก์เมตร มีปริมาตรเท่ากับกี่ลิตร?
\begin{enumerate}
\item 100 L
\item 1,000 L
\item 10,000 L
\item 1,000,000 L
\end{enumerate}
\textbf{คำตอบ: } ข
\textit{เฉลย: } 1 m³ = 1,000 L เพราะ 1 L = 1,000 cm³ และ 1 m³ = 1,000,000 cm³ ดังนั้น 1 m³ = 1,000 L
\HRule
\item[10] นักเรียนใช้ไม้บรรทัดวัดความยาวดินสอได้ $15.2$ cm ถ้าความไม่แน่นอนของไม้บรรทัดคือ $\pm 0.1$ cm จงเขียนผลการวัดในรูปแบบที่ถูกต้อง
\textbf{คำตอบ: } $15.2 \pm 0.1$ cm
\textit{วิธีทำ: } เมื่อเขียนผลการวัดที่มีความไม่แน่นอน ต้องระบุค่าที่วัดได้ตามด้วย $\pm$ ความไม่แน่นอน
ดังนั้นผลการวัดคือ $15.2 \pm 0.1$ cm
\HRule
\item[11] ระยะทาง $1$ เมตรในระบบ SI ถูกนิยามจากข้อใด?
\begin{enumerate}
\item ระยะทางที่แสงเดินทางในสุญญากาศในเวลา $\frac{1}{3 \times 10^8}$ วินาที
\item ความยาวของแท่งพลาตินัม-อิริเดียมที่เก็บไว้ที่ $0^\circ$C
\item ระยะทางที่แสงเดินทางในสุญญากาศในเวลา $\frac{1}{299,792,458}$ วินาที
\item ระยะทางเท่ากับความยาวคลื่นของแสงสีแดง
\end{enumerate}
\textbf{คำตอบ: } ค
\textit{เฉลย: } นิยามปัจจุบันของเมตรคือ ระยะทางที่แสงเดินทางในสุญญากาศในเวลา $\frac{1}{299,792,458}$ วินาที (ซึ่งเท่ากับ $\frac{1}{c}$ วินาที โดย c = ความเร็วแสง)
\HRule
\item[12] ข้อใดเรียงลำดับจากมากไปน้อยได้ถูกต้อง?
\begin{enumerate}
\item G > M > k > m
\item T > G > M > m
\item k > M > G > T
\item m > \mu > n > p
\end{enumerate}
\textbf{คำตอบ: } ก
\textit{เฉลย: } ค่าของคำอุปสรรค: G (giga) = $10^9$, M (mega) = $10^6$, k (kilo) = $10^3$, m (milli) = $10^{-3}$
ดังนั้นลำดับจากมากไปน้อยคือ G > M > k > m
\HRule
\item[13] ที่ดินแปลงหนึ่งมีพื้นที่ $2.5$ ไร่ จะมีพื้นที่กี่ตารางเมตร?
(กำหนดให้ 1 ไร่ = 400 ตารางวา, 1 ตารางวา = 4 ตารางเมตร)
\textbf{คำตอบ: } 4000 ตารางเมตร
\textit{วิธีทำ: } $2.5$ ไร่ = $2.5 \times 400$ ตารางวา = $1,000$ ตารางวา
$1,000$ ตารางวา = $1,000 \times 4$ ตารางเมตร = $4,000$ ตารางเมตร
\HRule
\item[14] ผลลัพธ์ของ $\frac{8.0}{2.00}$ มีเลขนัยสำคัญกี่ตัว?
\begin{enumerate}
\item 1 ตัว
\item 2 ตัว
\item 3 ตัว
\item 4 ตัว
\end{enumerate}
\textbf{คำตอบ: } ข
\textit{เฉลย: } $8.0$ มี 2 ตัว, $2.00$ มี 3 ตัว คำตอบต้องมีเลขนัยสำคัญน้อยที่สุด คือ 2 ตัว
$\frac{8.0}{2.00} = 4.0$ (มี 2 ตัว)
\HRule
\item[15] นักวิ่งวิ่งด้วยความเร็ว $36$ km/h ถ้าเขาวิ่งนาน $10$ นาที จะวิ่งได้ระยะทางกี่เมตร?
\begin{enumerate}
\item 360 m
\item 600 m
\item 3600 m
\item 6000 m
\end{enumerate}
\textbf{คำตอบ: } ง
\textit{เฉลย: } $36$ km/h = $36 \times \frac{1000}{3600}$ m/s = $10$ m/s
ระยะทาง = ความเร็ว $\times$ เวลา = $10 \times (10 \times 60)$ = $10 \times 600$ = $6,000$ m
\HRule
\item[16] นักเรียนชั่งน้ำหนักได้ $55.0$ kg และวัดส่วนสูงได้ $165.0$ cm จงหาค่า BMI (ดัชนีมวลร่างกาย) ของนักเรียน
(สูตร: BMI = มวล/ส่วนสูง² โดยมวลมีหน่วย kg และส่วนสูงมีหน่วยเมตร)
\textbf{คำตอบ: } 20.2 kg/m²
\textit{วิธีทำ: } ส่วนสูง = $165.0$ cm = $1.650$ m
BMI = $\frac{55.0}{(1.650)^2} = \frac{55.0}{2.7225} = 20.20$
มวลมี 3 ตัว, ส่วนสูงมี 4 ตัว ดังนั้นตอบ 3 ตัว: $20.2$ kg/m²
\HRule
\item[17] เมื่อวัดความยาวดินสอด้วยไม้บรรทัดที่มีสเกลละเอียดถึง mm พบว่าความยาวอยู่ระหว่าง $15.2$ cm กับ $15.3$ cm ความไม่แน่นอนในการวัดควรเป็นเท่าไหร่?
\begin{enumerate}
\item $\pm 0.05$ cm
\item $\pm 0.1$ cm
\item $\pm 0.5$ cm
\item $\pm 1.0$ cm
\end{enumerate}
\textbf{คำตอบ: } ก
\textit{เฉลย: } ความไม่แน่นอนของการวัด = ครึ่งหนึ่งของความละเอียดของเครื่องมือ
ความละเอียดของไม้บรรทัด = $1$ mm = $0.1$ cm
ความไม่แน่นอน = $\pm \frac{0.1}{2} = \pm 0.05$ cm
\HRule
\item[18] จูล (J) เป็นหน่วยอนุพัทธ์ของพลังงาน ซึ่งสามารถเขียนในรูปหน่วยฐาน SI ได้ตรงกับข้อใด?
\begin{enumerate}
\item kg·m·s⁻²
\item kg·m²·s⁻²
\item kg·m²·s⁻¹
\item kg·m·s⁻¹
\end{enumerate}
\textbf{คำตอบ: } ข
\textit{เฉลย: } พลังงาน = งาน = แรง × ระยะทาง = (kg·m·s⁻²) × m = kg·m²·s⁻²
ดังนั้น J = kg·m²·s⁻²
\HRule
\item[19] ถังน้ำทรงสี่เหลี่ยมมีขนาด กว้าง $50$ cm ยาว $80$ cm สูง $40$ cm ถังนี้จุน้ำได้กี่ลิตร?
\textbf{คำตอบ: } 160 L
\textit{วิธีทำ: } ปริมาตร = กว้าง × ยาว × สูง = $50 \times 80 \times 40$ cm³ = $160,000$ cm³
เพราะ $1$ L = $1,000$ cm³ ดังนั้น ถังจุน้ำได้ $\frac{160,000}{1000} = 160$ L
\HRule
\item[20] ถ้าให้ $g = 9.81$ m/s² และต้องการคำนวณโดยใช้เลขนัยสำคัญ 4 ตัวเท่านั้น ค่าที่ใช้ควรเป็นเท่าไหร่?
\begin{enumerate}
\item 9.810 m/s²
\item 9.8 m/s²
\item 9.811 m/s²
\item 9.8100 m/s²
\end{enumerate}
\textbf{คำตอบ: } ก
\textit{เฉลย: } $g = 9.81$ m/s² มีเลขนัยสำคัญ 3 ตัวอยู่แล้ว
ต้องการ 4 ตัว ต้องเขียนเป็น $9.810$ (เติมศูนย์ท้ายเพื่อบอกความละเอียด)
หมายเหตุ: $9.81$ กับ $9.810$ มีค่าเท่ากัน แต่ $9.810$ บอกว่าวัดได้ละเอียดถึง 3 ตำแหน่งทศนิยม
\HRule
\item[21] วัดมวลได้ $100.0 \pm 0.1$ g และปริมาตรได้ $50.0 \pm 0.5$ cm³ จงหาความหนาแน่นพร้อมความไม่แน่นอน
(ความหนาแน่น = มวล/ปริมาตร)
\textbf{คำตอบ: } $2.00 \pm 0.03$ g/cm³
\textit{วิธีทำ: } ความหนาแน่น = $\frac{100.0}{50.0} = 2.00$ g/cm³
ความไม่แน่นอนสัมพัทธ์: $\frac{\Delta m}{m} + \frac{\Delta V}{V} = \frac{0.1}{100.0} + \frac{0.5}{50.0} = 0.001 + 0.01 = 0.011$
$\Delta \rho = 2.00 \times 0.011 = 0.022 \approx 0.03$ g/cm³
ดังนั้น $\rho = 2.00 \pm 0.03$ g/cm³
\HRule
\item[22] นิยามของ 1 วินาทีในระบบ SI คือข้อใด?
\begin{enumerate}
\item ระยะเวลาที่ประจุไฟฟ้าเคลื่อนที่ผ่านความต้านทาน 1 โอห์ม
\item ระยะเวลาที่อะตอมซีเซียม-133 สั่น 9,192,631,770 รอบ
\item ระยะเวลาที่โลกหมุนรอบตัวเอง 1 รอบ
\item ระยะเวลาที่แสงเดินทาง 299,792,458 เมตร
\end{enumerate}
\textbf{คำตอบ: } ข
\textit{เฉลย: } 1 วินาที = ระยะเวลาที่อะตอมซีเซียม-133 สั่น $9,192,631,770$ รอบ (ใช้การสั่นของอะตอมซึ่งมีความแม่นยำสูงมาก)
\HRule
\item[23] ความหนาแน่นของน้ำมันเท่ากับ $0.80$ g/cm³ ความหนาแน่นนี้มีค่ากี่ kg/m³?
\textbf{คำตอบ: } 800 kg/m³
\textit{วิธีทำ: } $0.80$ g/cm³ = $0.80 \times \frac{10^{-3}}{10^{-6}}$ kg/m³ = $0.80 \times 10^3$ kg/m³ = $800$ kg/m³
\HRule
\item[24] นักเรียนวัดเวลาการตกของลูกกอล์ฟ 5 ครั้ง ได้ผลดังนี้ 2.31, 2.33, 2.30, 2.32, 2.31 วินาที ค่าเฉลี่ยและความไม่แน่นอนของการวัดเป็นเท่าไหร่?
\begin{enumerate}
\item $2.31 \pm 0.01$ s
\item $2.31 \pm 0.02$ s
\item $2.32 \pm 0.01$ s
\item $2.31 \pm 0.05$ s
\end{enumerate}
\textbf{คำตอบ: } ก
\textit{เฉลย: } ค่าเฉลี่ย = $\frac{2.31+2.33+2.30+2.32+2.31}{5} = \frac{11.57}{5} = 2.314 \approx 2.31$ s
ส่วนเบี่ยงเบนมาตรฐาน = $\sqrt{\frac{\sum(x_i - \bar{x})^2}{n-1}} = \sqrt{\frac{0.0006}{4}} \approx 0.012$ s
ปัดเศษความไม่แน่นอน = $\pm 0.01$ s
ดังนั้น $t = 2.31 \pm 0.01$ s
\HRule
\item[25] จงพิจารณาข้อความต่อไปนี้:
ก. ค่า $g = 9.8$ m/s² มีเลขนัยสำคัญ 2 ตัว
ข. การบวก $12.57 + 3.7 = 16.3$
ค. การคูณ $4.56 \times 2.0 = 9.1$
ข้อใดถูกต้อง?
\begin{enumerate}
\item ก และ ข
\item ก และ ค
\item ข และ ค
\item ก, ข และ ค
\end{enumerate}
\textbf{คำตอบ: } ง
\textit{เฉลย: } ก: $9.8$ มี 2 ตัว ✓
ข: $12.57 + 3.7$ ตอบ $16.3$ (ตำแหน่งทศนิยมน้อยสุดคือ 1 ตำแหน่ง จาก 3.7) ✓
ค: $4.56 \times 2.0 = 9.12 \approx 9.1$ (เลขนัยสำคัญน้อยสุดคือ 2 ตัว จาก 2.0) ✓
ดังนั้นข้อ ง ถูกต้อง
\HRule
\end{enumerate}
\newpage
\section*{แบบฝึกเพิ่มเติม (30 ข้อ)}
\begin{enumerate}
\item[1] ข้อใดไม่ใช่ปริมาณมูลฐานในระบบ SI?
\begin{enumerate}
\item มวล
\item น้ำหนัก
\item กระแสไฟฟ้า
\item อุณหภูมิ
\end{enumerate}
\textbf{คำตอบ: } ข
\textit{เฉลย: } น้ำหนักเป็นปริมาณอนุพัทธ์ (น้ำหนัก = มวล × g) ส่วนมวล กระแสไฟฟ้า และอุณหภูมิเป็นปริมาณมูลฐาน
\HRule
\item[2] $250$ เซนติเมตร เท่ากับกี่เมตร?
\begin{enumerate}
\item 2.5 m
\item 25 m
\item 0.25 m
\item 2500 m
\end{enumerate}
\textbf{คำตอบ: } ก
\textit{เฉลย: } $250$ cm = $250 \times 10^{-2}$ m = $2.5$ m
\HRule
\item[3] $0.000001$ เมตร มีค่าเท่ากับเท่าไหร่ในหน่วยไมโครเมตร?
\begin{enumerate}
\item 1 $\mu$m
\item 10 $\mu$m
\item 100 $\mu$m
\item 0.1 $\mu$m
\end{enumerate}
\textbf{คำตอบ: } ก
\textit{เฉลย: } $0.000001$ m = $10^{-6}$ m = $1$ $\mu$m (ไมโครเมตร)
\HRule
\item[4] จำนวน $700$ มีเลขนัยสำคัญกี่ตัว?
\begin{enumerate}
\item 1 ตัว
\item 2 ตัว
\item 3 ตัว
\item ขึ้นกับบริบท
\end{enumerate}
\textbf{คำตอบ: } ง
\textit{เฉลย: } 700 อาจมี 1 ตัว (7) หรือ 2 ตัว (7 และ 0 ท้าย ถ้ามีจุดทศนิยม) หรือ 3 ตัว (ถ้าเขียน 700.) ต้องดูบริบทการใช้งาน
\HRule
\item[5] จงแปลง $3.5$ กิโลกรัมเป็นกรัม
\textbf{คำตอบ: } 3500 g
\textit{วิธีทำ: } $3.5$ kg = $3.5 \times 1000$ g = $3500$ g
\HRule
\item[6] ข้อใดเรียงลำดับจากน้อยไปมากได้ถูกต้อง?
\begin{enumerate}
\item nm < $\mu$m < mm < cm
\item nm > $\mu$m > mm > cm
\item nm < mm < $\mu$m < cm
\item nm > mm > $\mu$m > cm
\end{enumerate}
\textbf{คำตอบ: } ก
\textit{เฉลย: } nm ($10^{-9}$) < $\mu$m ($10^{-6}$) < mm ($10^{-3}$) < cm ($10^{-2}$)
\HRule
\item[7] อุณหภูมิ $27^\circ$C เท่ากับกี่เคลวิน?
\begin{enumerate}
\item 246 K
\item 273 K
\item 300 K
\item 246 K
\end{enumerate}
\textbf{คำตอบ: } ค
\textit{เฉลย: } $T(K) = T(^\circ C) + 273.15$ ดังนั้น $27 + 273.15 = 300.15 \approx 300$ K
\HRule
\item[8] จงบวก $23.6 + 8.27 + 3.451$ โดยใช้เลขนัยสำคัญที่ถูกต้อง
\textbf{คำตอบ: } 35.3
\textit{วิธีทำ: } ตำแหน่งทศนิยมน้อยสุดคือ 1 ตำแหน่ง (จาก 23.6)
$23.6 + 8.27 + 3.451 = 35.321 \approx 35.3$
\HRule
\item[9] ผลลัพธ์ของ $5 \times 4.2$ มีเลขนัยสำคัญกี่ตัว?
\begin{enumerate}
\item 1 ตัว
\item 2 ตัว
\item 3 ตัว
\item 4 ตัว
\end{enumerate}
\textbf{คำตอบ: } ข
\textit{เฉลย: } 5 มี 1 ตัว, 4.2 มี 2 ตัว ดังนั้นผลลัพธ์ต้องมี 1 ตัว: $5 \times 4.2 = 21 \approx 20$
\HRule
\item[10] รถจักรยานยนต์วิ่งด้วยความเร็ว $54$ km/h จงหาความเร็วนี้ในหน่วย m/s
\textbf{คำตอบ: } 15 m/s
\textit{วิธีทำ: } $54$ km/h = $54 \times \frac{5}{18}$ m/s = $15$ m/s
\HRule
\item[11] นิวตัน (N) สามารถเขียนในรูปหน่วยฐาน SI ได้ตรงกับข้อใด?
\begin{enumerate}
\item kg·m·s⁻¹
\item kg·m·s⁻²
\item kg·m²·s⁻¹
\item kg·m²·s⁻²
\end{enumerate}
\textbf{คำตอบ: } ข
\textit{เฉลย: } แรง = มวล × ความเร่ง = kg × m/s² = kg·m·s⁻²
\HRule
\item[12] เครื่องชั่งที่มีความละเอียด $0.01$ kg วัดมวลได้ $5.67$ kg ผลการวัดควรเขียนอย่างไร?
\begin{enumerate}
\item $5.67 \pm 0.01$ kg
\item $5.67 \pm 0.005$ kg
\item $5.7 \pm 0.1$ kg
\item $5.67 \pm 0.5$ kg
\end{enumerate}
\textbf{คำตอบ: } ก
\textit{เฉลย: } ความไม่แน่นอน = $\pm \frac{1}{2}$ ของความละเอียด = $\pm 0.005$ kg ปัดเป็น 1 ตำแหน่ง = $\pm 0.01$ kg
\HRule
\item[13] สนามรูปสี่เหลี่ยมผืนผ้ากว้าง $40$ m ยาว $60$ m มีพื้นที่กี่ตารางกิโลเมตร?
\textbf{คำตอบ: } $2.4 \times 10^{-6}$ km²
\textit{วิธีทำ: } พื้นที่ = $40 \times 60$ = $2400$ m²
$2400$ m² = $2400 \times 10^{-6}$ km² = $2.4 \times 10^{-3}$ km²
ผิด: $2400$ m² = $2400 \times (10^{-3})^2$ km² = $2400 \times 10^{-6}$ km² = $2.4 \times 10^{-3}$ km²
\HRule
\item[14] ข้อใดมีเลขนัยสำคัญ 4 ตัว?
\begin{enumerate}
\item 0.0400
\item 4000
\item 4.000
\item 4.0 \times 10^3
\end{enumerate}
\textbf{คำตอบ: } ค
\textit{เฉลย: } ก มี 3 ตัว (4, 0, 0), ข มี 1 ตัว, ง มี 2 ตัว, ค 4.000 มี 4 ตัว (4, 0, 0, 0)
\HRule
\item[15] จำนวนวินาทีใน 1 วัน (24 ชั่วโมง) มีค่าเท่าไหร่?
\textbf{คำตอบ: } 86400 s
\textit{วิธีทำ: } 1 วัน = 24 ชม. = $24 \times 60$ นาที = $1440$ นาที
$= 1440 \times 60$ วินาที = $86,400$ วินาที
\HRule
\item[16] นกบินด้วยความเร็ว $30$ m/s ความเร็วนี้เท่ากับกี่ km/h?
\begin{enumerate}
\item 90 km/h
\item 108 km/h
\item 120 km/h
\item 150 km/h
\end{enumerate}
\textbf{คำตอบ: } ข
\textit{เฉลย: } $30$ m/s = $30 \times \frac{18}{5}$ km/h = $108$ km/h
\HRule
\item[17] จงหาผลลัพธ์ของ $\frac{9.84}{2.4}$ โดยใช้เลขนัยสำคัญที่ถูกต้อง
\textbf{คำตอบ: } 4.1
\textit{วิธีทำ: } 9.84 มี 3 ตัว, 2.4 มี 2 ตัว ดังนั้นตอบ 2 ตัว
$\frac{9.84}{2.4} = 4.1$
\HRule
\item[18] มาตรฐานของหน่วยกิโลกรัมในปัจจุบัน (2019) ถูกนิยามจากข้อใด?
\begin{enumerate}
\item มวลของแท่งพลาตินัม-อิริเดียมที่เก็บไว้ในกรุงปารีส
\item มวลของอะตอมคริปตอน-84 จำนวนหนึ่ง
\item ค่าคงที่พลังค์ที่กำหนดให้มีค่าตายตัว
\item มวลของน้ำบริสุทธิ์ 1 ลิตรที่ $4^\circ$C
\end{enumerate}
\textbf{คำตอบ: } ค
\textit{เฉลย: } ตั้งแต่ปี 2019 กิโลกรัมถูกนิยามจากค่าคงที่พลังค์ (h) ที่กำหนดให้มีค่า $6.62607015 \times 10^{-34}$ kg·m²·s⁻¹ แทนการใช้แท่งพลาตินัม-อิริเดียม
\HRule
\item[19] สารมีความหนาแน่น $2.7$ g/cm³ ความหนาแน่นนี้เท่ากับกี่ kg/m³?
\textbf{คำตอบ: } 2700 kg/m³
\textit{วิธีทำ: } $2.7$ g/cm³ = $2.7 \times \frac{10^{-3}}{10^{-6}}$ kg/m³ = $2.7 \times 10^3$ kg/m³ = $2700$ kg/m³
\HRule
\item[20] ค่าคงตัว $c = 3.00 \times 10^8$ m/s มีเลขนัยสำคัญกี่ตัว?
\begin{enumerate}
\item 1 ตัว
\item 2 ตัว
\item 3 ตัว
\item 9 ตัว
\end{enumerate}
\textbf{คำตอบ: } ค
\textit{เฉลย: } 3.00 มีเลขนัยสำคัญ 3 ตัว (3, 0, 0) ศูนย์ท้ายในทศนิยมมีนัยสำคัญ
\HRule
\item[21] แรง $500$ นิวตัน มีค่าเท่ากับกี่กิโลนิวตัน?
\textbf{คำตอบ: } 0.5 kN
\textit{วิธีทำ: } $500$ N = $500 \times 10^{-3}$ kN = $0.5$ kN
\HRule
\item[22] วัดความยาว 4 ครั้งได้ 10.2, 10.4, 10.1, 10.3 cm ค่าเฉลี่ยและความไม่แน่นอนคือข้อใด?
\begin{enumerate}
\item $10.25 \pm 0.13$ cm
\item $10.25 \pm 0.06$ cm
\item $10.3 \pm 0.1$ cm
\item $10.2 \pm 0.1$ cm
\end{enumerate}
\textbf{คำตอบ: } ค
\textit{เฉลย: } ค่าเฉลี่ย = $\frac{10.2+10.4+10.1+10.3}{4} = 10.25 \approx 10.3$
ส่วนเบี่ยงเบนมาตรฐาน = $\sqrt{\frac{0.05^2+0.15^2+0.15^2+0.05^2}{3}} \approx 0.13$
ปัดเป็น 1 ตัวเลขนัยสำคัญ = $\pm 0.1$ cm
\HRule
\item[23] ขวดน้ำมีปริมาตร $1500$ cm³ จงเขียนในรูปลิตร
\textbf{คำตอบ: } 1.5 L
\textit{วิธีทำ: } $1500$ cm³ = $1500 \times 10^{-3}$ L = $1.5$ L (เพราะ $1$ L = $1000$ cm³)
\HRule
\item[24] จงหาผลลัพธ์ของ $(2.36)^2$ โดยใช้เลขนัยสำคัญที่ถูกต้อง
\begin{enumerate}
\item 5.57
\item 5.570
\item 5.6
\item 5.5696
\end{enumerate}
\textbf{คำตอบ: } ค
\textit{เฉลย: } $2.36$ มี 3 ตัว ดังนั้นผลลัพธ์ต้องมี 3 ตัว
$(2.36)^2 = 5.5696 \approx 5.57$
แต่เมื่อปัดเศษตามการคูณหาร (3 ตัว) ได้ 5.57 ซึ่งมี 3 ตัว
\HRule
\item[25] เครื่องปรับอากาศมีกำลัง $1.5$ กิโลวัตต์ มีกำลังเท่ากับกี่วัตต์?
\textbf{คำตอบ: } 1500 W
\textit{วิธีทำ: } $1.5$ kW = $1.5 \times 1000$ W = $1500$ W
\HRule
\item[26] $10^{-9}$ เรียกว่าอะไร?
\begin{enumerate}
\item มิลลิ
\item ไมโคร
\item นาโน
\item พิโก
\end{enumerate}
\textbf{คำตอบ: } ค
\textit{เฉลย: } $10^{-9}$ = นาโน (nano, n)
\HRule
\item[27] อาหารจานหนึ่งให้พลังงาน $500$ กิโลแคลอรี จงเปลี่ยนเป็นจูล
(กำหนดให้ 1 cal = 4.184 J)
\textbf{คำตอบ: } 2092 J
\textit{วิธีทำ: } $500$ kcal = $500 \times 1000$ cal = $500,000$ cal
$= 500,000 \times 4.184$ J = $2,092,000$ J = $2.092 \times 10^6$ J
ปัดเป็น $2.1 \times 10^6$ J หรือ $2,092,000$ J
\HRule
\item[28] นักเรียนวัดเวลาการแกว่งของลูกตุ้ม 5 รอบ ได้ $12.35$ วินาที ถ้าความไม่แน่นอนของนาฬิกาจับเวลาคือ $\pm 0.2$ วินาที คาบการแกว่ง (ต่อรอบ) มีค่าเท่าไหร่?
\begin{enumerate}
\item $2.47 \pm 0.04$ s
\item $2.47 \pm 0.02$ s
\item $2.5 \pm 0.2$ s
\item $2.47 \pm 0.2$ s
\end{enumerate}
\textbf{คำตอบ: } ก
\textit{เฉลย: } คาบ = $\frac{12.35}{5} = 2.47$ s
ความไม่แน่นอนของคาบ = $\frac{0.2}{5} = 0.04$ s
ดังนั้น $T = 2.47 \pm 0.04$ s
\HRule
\item[29] จงหาค่า $\sqrt{144}$ โดยใช้เลขนัยสำคัญที่ถูกต้อง
\textbf{คำตอบ: } 12.0
\textit{วิธีทำ: } $\sqrt{144} = 12$ มีเลขนัยสำคัญ 1 ตัว หรือ 3 ตัวถ้าเขียน 12.0
ปกติ 144 ถือว่ามี 3 ตัว (ตัวเลขทั้งหมดไม่ใช่ศูนย์นำหน้า)
ดังนั้น $\sqrt{144} = 12.0$ (3 ตัว) หรือ 12 (1 ตัว) ขึ้นกับบริบท
\HRule
\item[30] ปริมาณน้ำฝนที่ตกในพื้นที่หนึ่งเป็น $25$ มิลลิเมตร หมายความว่าอย่างไร?
\begin{enumerate}
\item ฝนตกหนักมาก
\item มีน้ำฝนสะสมสูง 25 มิลลิเมตรบนพื้นที่นั้น
\item มีฝนตกเป็นเส้นยาว 25 มิลลิเมตร
\item ถูกทั้งข้อ ก และ ข
\end{enumerate}
\textbf{คำตอบ: } ข
\textit{เฉลย: } ปริมาณน้ำฝนวัดเป็น mm หมายถึงความสูงของน้ำที่สะสมบนพื้นผิวเมื่อฝนหยุดตก ถ้าฝนตก 25 mm หมายความว่ามีน้ำสูง 25 mm บนพื้นที่นั้น
\HRule
\end{enumerate}
\end{document}
